% Standalone version. The manuscript-insertable body is v2_methods_section.tex.
% This file embeds that body so the Codex LaTeX compiler needs no project files.
\documentclass[10pt]{article}
\usepackage[margin=0.9in]{geometry}
\usepackage[T1]{fontenc}
\usepackage{lmodern}
\usepackage{amsmath,amssymb}
\usepackage{booktabs,tabularx}
\usepackage{algorithm}
\usepackage{algpseudocode}
\usepackage[hidelinks]{hyperref}
\setlength{\emergencystretch}{2em}
\begin{document}
% BEGIN EMBEDDED METHODS SECTION
% Manuscript-ready section; include in a document that loads:
% amsmath, amssymb, booktabs, tabularx, algorithm, algpseudocode.
% Audited against branch value-aware_cross-step-reuse_v2, commit 37bfbae3.
% Main sources: v31_value_selectors.py:137--183;
% v31_value_kernels.py:18--148,409--439;
% vllm_adapter.py:1110--1257,1395--1406,1822--1879;
% diffusion_gemma_jl_output_aware/projections.py:21--46;
% v31_value_projection.py:14--54; v27_fa4.py:98--134.
% Native model: pinned revision f7f5b7f5fa82ffc52addd066915886d497f5517b.
% IMPORTANT: tau_V=0.021 instantiates the recent evaluated configuration;
% it is not an intrinsic constant or a newly selected/calibrated threshold.
% Earlier tau_V=0.027 and 0.015 results are different configurations.
% Provenance: Yuhan introduced the inherited V31 cross-step reuse lifecycle.
% Gaussian32/value-direction machinery originates in Junyu Liao's prior work.
% The present V2 selector integration is this study's separately named variant.
% Exact projection seed material, for each native KV head g:
% jl_output_v1/gaussian/32/1729/{layer}/{head}/{width}
% derived = int.from_bytes(sha256(material.encode()).digest()[:8], 'little')
%           % (2**63-1); private CPU torch.Generator, torch.randn / sqrt(32).

\section{Online Value-Direction-Aware Block Selection with Cross-Step Reuse}
\label{sec:v2-method}

Our method, V2, constructs a block-sparse attention map using an ordered
value-direction criterion and reuses that map across denoising calls.
The criterion estimates each candidate block's effect on a running projected
attention output. Its defining property is that skipped blocks advance the
running mass while leaving the normalized routing output unchanged.
We combine this selector with the inherited V31 map-reuse lifecycle.\footnote{The
cross-step reuse lifecycle originates in Yuhan's V31 implementation; the
Gaussian32 machinery follows Junyu Liao's prior value-direction implementation.
V2 denotes the selector integration studied here.}

\subsection{Attention scope and block statistics}
\label{sec:v2-statistics}

We apply sparsity only to the denoising calls of the five GLOBAL layers of
DiffusionGemma-26B-A4B. Each layer has $H_q=16$ query heads and $H_{kv}=2$
native KV heads, with $d_k=d_v=512$. Each query head independently selects
KV blocks for every query block of $B_q=128$ rows; each KV block contains
$B_k=64$ keys. A canvas has at most $C=256$ query rows. LOCAL layers,
prompt prefill, and committed-canvas encoding retain their native behavior
and execute without this sparsity policy.

At a selection event, suppress the layer and query-head indices. Let
$\mathcal Q_b$ be the valid rows of query block $b$, $P$ the prefix length,
and $\mathcal B_j$ the $j$th KV block in increasing token order. The keys
consist of the prefix followed by the current canvas. We use current,
native post-normalization and post-positional-encoding queries and keys:
\begin{equation}
 s_{iu}=\alpha q_i^{\mathsf T}k_u,
 \qquad z_u=R_{\ell,g(h)}^{\mathsf T}v_u,
 \qquad R_{\ell,g(h)}\in\mathbb R^{d_v\times r},
 \quad R_{ab}\sim\mathcal N(0,1/r).
 \label{eq:v2-projection}
\end{equation}
Here $g(h)=\lfloor h/(H_q/H_{kv})\rfloor$ is the native GQA mapping and
$r=32$. The native scale is $\alpha=1$ in the pinned Gemma4 implementation.
Projection matrices are fixed per layer and native KV head, using base seed
$\sigma_R=1729$ and deterministic SHA256-derived seeds in private CPU random
generators. Matrices are reused, whereas current values are reprojected at
every selection event; projection generation does not consume sampler RNG.

For each row and block, compute
\begin{align}
 Z_{ij}&=\sum_{u\in\mathcal B_j\cap\mathcal K_i}\exp(s_{iu}),
 & L_{ij}&=\log Z_{ij}, \label{eq:v2-mass}\\
 \mu_{ij}&=\frac{1}{Z_{ij}}
       \sum_{u\in\mathcal B_j\cap\mathcal K_i}\exp(s_{iu})z_u,
 & \nu_{\ell g}&=\max\!\left\{\epsilon,
       \sqrt{\frac{1}{|\mathcal K|}
       \sum_{u\in\mathcal K}\|v_u\|_2^2}\right\}.
 \label{eq:v2-mean-scale}
\end{align}
$\mathcal K_i$ contains structurally valid keys for row $i$; GLOBAL
denoising attention is bidirectional over the prefix and current canvas.
$\mathcal K$ contains all valid current keys, including canvas keys.
The scale $\nu_i=\nu_{\ell,g(h)}$ uses full-dimensional values and is shared
within a native KV head. For empty row/block intersections, set
$L_{ij}=-\infty$ and $\mu_{ij}=0$. All reductions are numerically stabilized.
The quantity $\mu_{ij}$ is an attention-weighted within-block mean.

\subsection{Mass-inclusive online selection}
\label{sec:v2-selection}

At each selection event, initialize the running mass $A_i=0$, normalized
projected routing output $\widehat o_i=0$, and retained-support indicator
$u_i=0$. Scan all KV blocks from oldest to newest, including the protected
canvas blocks. Before block $j$, calculate
\begin{equation}
 \eta_{ij}=\frac{Z_{ij}}{A_i+Z_{ij}},
 \qquad
 \rho_{bj}=\max_{i\in\mathcal I_{bj}}
 \frac{\eta_{ij}\|\mu_{ij}-\widehat o_i\|_2}{\nu_i},
 \quad
 \mathcal I_{bj}=\{i\in\mathcal Q_b:L_{ij}>-\infty\}.
 \label{eq:v2-risk}
\end{equation}
Set $\eta_{ij}=0$ outside $\mathcal I_{bj}$ and $\rho_{bj}=0$ if this set
is empty. The vector $\eta_{ij}(\mu_{ij}-\widehat o_i)$ is the hypothetical
change in the routing output when incorporating the block. The maximum
permits any valid row to retain the entire physical tile.

For an unprotected block with existing retained support, retain it if
\begin{equation}
 \log\rho_{bj}+\lambda\overline M_{bj}\ \ge\ \log\tau_V.
 \label{eq:v2-threshold}
\end{equation}
$\tau_V\ge0$ is one shared threshold across all GLOBAL layers, heads,
query blocks, and selection events in an experimental configuration.
$\overline M_{bj}$ indicates membership in the currently held map at a
within-canvas refresh; it is zero at each canvas's initial selection.
The sticky bonus $\lambda=1.386$ lowers a held tile's effective retention
threshold by a factor $\exp(\lambda)\approx4$. Equality retains the tile.
We additionally retain every tile intersecting the current canvas
($j\ge\lfloor P/B_k\rfloor$), including a mixed prefix/canvas boundary
tile, and every tile providing the first retained support for any active
row. Invalid logits or projected means cause all tiles of the affected
query-head/query-block unit to be retained. Padded and empty rows do not vote.

With retained-tile indicator $M_{bj}\in\{0,1\}$, update
\begin{align}
 A_i'&=A_i+Z_{ij}, \label{eq:v2-mass-update}\\
 \widehat o_i'&=
 \begin{cases}
  (1-\eta_{ij})\widehat o_i+\eta_{ij}\mu_{ij}, & M_{bj}=1,\\
  \widehat o_i, & M_{bj}=0.
 \end{cases}
 \label{eq:v2-output-update}
\end{align}
Thus $A_i$ includes every previously scanned block, including skipped
blocks. Equivalently, a skipped block contributes the virtual numerator
$Z_{ij}\widehat o_i$, preserving the normalized state. Online denotes this
ordered scan; its states restart at each selection event. Selection is
order dependent and has variable support, with no fixed top-$k$ quota.
With the convention $\log0=-\infty$, $\tau_V=0$ retains every active tile.

\begin{algorithm}[t]
\caption{V2 selection for one layer, query head, and query block}
\label{alg:v2-selector}
\small
\begin{algorithmic}[1]
\Require Current statistics $L_{ij},\mu_{ij},\nu_i$; prefix length $P$;
held mask $\overline M$ (zero at initial selection); $\tau_V,\lambda$
\Ensure Retained-block mask $M\in\{0,1\}^{J}$
\If{invalid logits or projected means affect this query unit}
    \State \Return $\boldsymbol 1_J$ \Comment{retain all tiles}
\EndIf
\State $\ell_i\gets-\infty$, $\widehat o_i\gets0$, $u_i\gets0$ for all $i\in\mathcal Q_b$
\State $M\gets\boldsymbol 0_J$
\For{$j=0,\ldots,J-1$ in increasing token order}
    \State $\mathcal I\gets\{i\in\mathcal Q_b:L_{ij}>-\infty\}$
    \State $\ell_i^+\gets\operatorname{logaddexp}(\ell_i,L_{ij})$ for all $i$
    \State $\eta_i\gets\exp(L_{ij}-\ell_i^+)$ for $i\in\mathcal I$; $0$ otherwise
    \State $v\gets\displaystyle\max_{i\in\mathcal I}
        \log\!\left(\eta_i\|\mu_{ij}-\widehat o_i\|_2/\nu_i\right)$
        \Comment{$\max\varnothing=-\infty$}
    \State $p\gets[j\ge\lfloor P/B_k\rfloor]$
    \State $f\gets[\exists i\in\mathcal I:u_i=0]$
    \State $M_j\gets[\mathcal I\ne\varnothing]\land
        \bigl(p\lor f\lor[v+\lambda\overline M_j\ge\log\tau_V]\bigr)$
    \If{$M_j=1$}
        \State $\widehat o_i\gets(1-\eta_i)\widehat o_i+\eta_i\mu_{ij}$ for $i\in\mathcal I$
        \State $u_i\gets1$ for $i\in\mathcal I$
    \EndIf
    \State $\ell_i\gets\ell_i^+$ for all $i$ \Comment{also on a skip}
\EndFor
\State \Return $M$
\end{algorithmic}
\end{algorithm}

\subsection{Map reuse, carry, and progress-triggered refresh}
\label{sec:v2-reuse}

Number denoising calls within a canvas from $t=0$. Call $0$ executes
full-support FA4 unless a valid previous-canvas map can be carried. Carry
requires an immediately preceding canvas, its previous total KV length
equal to the new prefix length, and compatible query-block and KV-tile
geometry. It preserves decisions only for tiles wholly inside the old
prefix; all newer tiles, including the previous and current canvases, are
retained. Call $1$ executes full-support FA4 and constructs a fresh V2 map,
which applies from the following call. Subsequent ordinary calls reuse the
held map. A progress refresh likewise produces a full-support output and
installs a replacement map for subsequent calls. The sticky bonus uses
only a map selected within the same canvas. Maps and request-specific
histories are cleared between requests.

After call $t$, let $a_{it}$ be the acceptance indicator obtained by
evaluating the native entropy-bound rule on the refresh hook's row statistics,
$y_{it}$ the argmax token of the temperature-scaled logits, and
$p_{it}^{\max}$ its probability. Track a re-noising exponential moving
average $d_{it}$ and an accepted stable-run length $r_{it}$:
\begin{align}
 d_{it}&=\gamma d_{i,t-1}+(1-\gamma)(1-a_{it}),
          \qquad d_{i,-1}=1, \label{eq:v2-renoise}\\
 r_{it}&=
 \begin{cases}
   a_{it}, & t=0,\\
   a_{it}(r_{i,t-1}+1), & t>0,\ y_{it}=y_{i,t-1},\\
   0, & \text{otherwise}.
 \end{cases}
 \label{eq:v2-run}
\end{align}
Define settledness and its mean over the $n_q$ valid canvas rows by
\begin{equation}
 c_{it}=\left(1-e^{-r_{it}/\tau_s}\right)(1-d_{it})
        \left(1-\sqrt{\max\{0,1-p_{it}^{\max}\}}\right),
 \qquad \phi_t=\frac{1}{n_q}\sum_i c_{it}.
 \label{eq:v2-progress}
\end{equation}
The first completed call satisfying $\phi_t\ge\delta$ schedules selection
at the next denoising call in every GLOBAL layer. We use $\gamma=0.65$,
$\tau_s=2.5$, and $\delta=0.15$, with one threshold crossing per canvas.
The threshold is latched after firing. A refresh is executed only if native
decoding requests the next call; a crossing after call $0$ coincides with
the regular call-$1$ selection. The clock resets per canvas and uses
absolute settledness. Its C-gate-derived statistics determine refresh timing;
V2's tile score uses unit row weights.

\begin{algorithm}[t]
\caption{V2 within the unchanged native denoising loop}
\label{alg:v2-reuse}
\small
\begin{algorithmic}[1]
\Require Native model and sampler; per-layer held maps with geometry metadata
\For{each canvas $c$ requested by native generation}
    \State Reset progress history; $\mathit{fired}\gets\mathsf{false}$;
           $\mathit{pending}\gets\mathsf{false}$; $t\gets0$
    \While{the native sampler requests another denoising forward}
        \State $\mathit{select}\gets(t=1)\lor\mathit{pending}$;
               $\mathit{pending}\gets\mathsf{false}$
        \For{each decoder layer $\ell$ in model order}
            \If{$\ell$ is LOCAL}
                \State Execute native LOCAL attention
            \ElsIf{$t=0$}
                \If{the previous-canvas map passes the carry checks}
                    \State Execute current-Q/K/V FA4 on the carried map
                \Else
                    \State Execute full-support FA4
                \EndIf
            \ElsIf{$\mathit{select}$}
                \State Execute full-support FA4; retain its output and prefix masses
                \State Compute current projections, tile means, and RMS scales
                \State $\overline M^\ell\gets M^\ell$ if selected in $c$; zero otherwise
                \State $M^\ell\gets$ Algorithm~\ref{alg:v2-selector} independently for all $(h,b)$
                \State Store $M^\ell$ and its current-canvas geometry metadata
            \Else
                \State Execute current-Q/K/V FA4 on held map $M^\ell$
            \EndIf
        \EndFor
        \State Execute the unchanged native sampler and stopping logic
        \If{$\neg\mathit{fired}$}
            \State Update Eqs.~\eqref{eq:v2-renoise}--\eqref{eq:v2-progress} from completed-call statistics
            \If{$\phi_t\ge\delta$}
                \State $\mathit{pending}\gets\mathsf{true}$;
                       $\mathit{fired}\gets\mathsf{true}$
            \EndIf
        \EndIf
        \State $t\gets t+1$
    \EndWhile
    \State Commit the canvas using the native encoder path
\EndFor
\end{algorithmic}
\end{algorithm}

\subsection{Executed attention and native decoding}
\label{sec:v2-execution}

On a held or carried call, FA4 evaluates attention using the current
queries, keys, and original full-dimensional values over the retained keys
$\mathcal S_{b,t}$:
\begin{equation}
 o_{it}=\frac{\sum_{u\in\mathcal S_{b,t}\cap\mathcal K_i}
                     \exp(s_{iu,t})v_{u,t}}
                  {\sum_{u\in\mathcal S_{b,t}\cap\mathcal K_i}
                     \exp(s_{iu,t})}.
 \label{eq:v2-executed}
\end{equation}
Only the block map is reused. Attention probabilities and outputs are
recomputed; current canvas K/V are refreshed each call. The projected
routing state is discarded after selection and never becomes model output.
The mass-inclusive recurrence is therefore a routing surrogate, whereas
Eq.~\eqref{eq:v2-executed} renormalizes over retained keys. The threshold
does not provide a bound on final dense-attention error or an exact physical
sparsity target.

Native proposal sampling, acceptance, re-noising, self-conditioning, and
adaptive stopping are preserved. With $D_{\max}=48$, the pinned runtime uses
\begin{equation}
 T_t=T_{\min}+(T_{\max}-T_{\min})
       \frac{\max\{D_{\max}-t,1\}}{D_{\max}},
 \qquad (T_{\min},T_{\max})=(0.4,0.8).
 \label{eq:v2-temperature}
\end{equation}
Acceptance follows the native entropy-bound rule: for row entropies sorted
as $H_{\pi(1),t}\le\cdots\le H_{\pi(C),t}$, position $\pi(k)$ is accepted
when $\sum_{m=1}^{k}H_{\pi(m),t}-\max_{m\le k}H_{\pi(m),t}\le B_{\rm ent}$,
with $B_{\rm ent}=0.1$. Native early stopping requires agreement of two
consecutive denoising argmax states and mean processed-logit entropy below
$\theta_{\rm conf}=0.005$; otherwise the $48$-call cap applies. The generation
configuration's stability parameter is $1$ (the runtime stores two states).
EOS handling follows the native model configuration; thinking mode and
output limits remain workload-specific. Selection and refresh execute
inside already requested model forwards. Their attention perturbations can
change the resulting native trajectory and call count through the unchanged
sampler.

\subsection{Implementation and hyperparameters}
\label{sec:v2-implementation}

Native Q/K/V tensors use BF16; projected values, tile summaries, and
selection states use FP32. Projection and projected-value products use
TF32x3 for FP32 inputs; the value-norm reduction preserves the inherited
norm-then-square rounding convention. The refresh hook obtains argmax,
top-token probability, and entropy from fused FP32 logit reductions without
replacing the native sampler's computations. Full-support FA4 writes the prefix tile log
masses without feeding the observer statistics into its output. A separate
current-QK/projected-value pass supplies $\mu_{ij}$; its prefix log masses
are replaced with the FA4-observed values. This implementation reprojects
all current values at each selection. The selector stores $\log A_i$ and
uses log-sum-exp updates, avoiding underflow relative to a future large
logit. GPU scans process $128$ KV tiles per chunk and carry the running
state across chunks, preserving the sequential decision rule. Retained maps
are converted to explicit FA4 block lists. Selection work is incurred only
at initial/refresh events. Physical consumer work includes full-support
bootstrap/selection calls and retained work on held/carried calls; the
additional discovery and selection passes are accounted for separately.

Table~\ref{tab:v2-hyperparameters} specifies the V2 configuration with
$\tau_V=0.021$. The threshold is fixed per experimental configuration and
shared across its initial and refresh selections. Other thresholds and
their selection procedures are reported as separate experimental settings.

\begin{table*}[t]
\centering
\caption{V2 hyperparameters, fixed model geometry, and execution settings.
Values specify the $\tau_V=0.021$ configuration. Workload-specific settings
are bound to the corresponding evaluation protocol.}
\label{tab:v2-hyperparameters}
\footnotesize
\setlength{\tabcolsep}{4pt}
\renewcommand{\arraystretch}{1.08}
\begin{tabularx}{\textwidth}{@{}l l X@{}}
\toprule
Parameter & Symbol & Value or policy \\
\midrule
\multicolumn{3}{@{}l}{\textit{Block selection}} \\
Uniform retention threshold & $\tau_V$ & $0.021$; shared across GLOBAL layers and selection events \\
Projection rank & $r$ & $32$ \\
Projection distribution & $R_{ab}$ & $\mathcal N(0,1/r)$; fixed per layer/native KV head \\
Projection base seed & $\sigma_R$ & $1729$; SHA256-derived per-layer/head CPU seeds \\
Reference-scale floor & $\epsilon$ & $10^{-12}$ \\
Sticky log-risk bonus & $\lambda$ & $1.386$; within-canvas refresh only \\
Row aggregation & --- & Maximum over active rows of each query block \\
Row sensitivity weights & $w_i$ & $1$ \\
Scan order & --- & Increasing KV token index, prefix then canvas \\
Fixed support quota & --- & None; threshold-based variable support \\
Mandatory canvas/boundary tiles & --- & All retained \\
Additional sink/recent protections & --- & $0$ tiles / $0$ tokens \\
Additional mass anchors & --- & $0$ tiles \\
\midrule
\multicolumn{3}{@{}l}{\textit{Reuse and progress clock}} \\
Initial selection call & $t_{\rm sel}$ & $1$ (zero-based); full-support output \\
Cross-canvas first-call carry & --- & Enabled, subject to continuation and geometry checks \\
Progress threshold & $\delta$ & $0.15$; one latched crossing per canvas \\
Re-noising EMA coefficient & $\gamma$ & $0.65$; initial state $d_{i,-1}=1$ \\
Stable-run time scale & $\tau_s$ & $2.5$; initial run length $0$ \\
Progress definition & $\phi_t$ & Absolute mean settledness; no relative rescaling \\
Periodic/cheap refresh & --- & Disabled; full-support progress refresh only \\
\midrule
\multicolumn{3}{@{}l}{\textit{Native model and sampler}} \\
Sparse GLOBAL layers & $\mathcal L_G$ & $\{5,11,17,23,29\}$, zero-based \\
Query / KV heads & $H_q,H_{kv}$ & $16,2$ in each GLOBAL layer \\
GLOBAL head dimensions & $d_k,d_v$ & $512,512$ \\
Query / KV tile sizes & $B_q,B_k$ & $128,64$ \\
Native attention scale & $\alpha$ & $1$ \\
Canvas length & $C$ & $256$; final valid length may be shorter \\
Maximum denoising calls & $D_{\max}$ & $48$ per canvas \\
Native temperature bounds & $T_{\min},T_{\max}$ & $0.4,0.8$; exact schedule in Eq.~\eqref{eq:v2-temperature} \\
Entropy-bound acceptance & $B_{\rm ent}$ & $0.1$ \\
Stopping entropy threshold & $\theta_{\rm conf}$ & $0.005$ \\
Native stability setting & --- & $1$; agreement of two consecutive argmax states \\
Thinking / output limit / generation seed & --- & Workload-specific; unchanged within each comparison \\
\midrule
\multicolumn{3}{@{}l}{\textit{Execution}} \\
Batch size / hardware & --- & $1$ / NVIDIA H100 (SM90) \\
Attention consumer & --- & Pinned FlashAttention-4 block-sparse interface \\
vLLM / decode graph mode & --- & $0.30.0$ / PIECEWISE \\
Native / routing precision & --- & BF16 / FP32 \\
Selector scan chunk & $J_{\rm chunk}$ & $128$ KV tiles; recurrence preserved across chunks \\
LOCAL / prefill / canvas commit & --- & Native dense execution \\
\bottomrule
\end{tabularx}
\vspace{2pt}
\begin{minipage}{\textwidth}
\footnotesize
The inherited launcher field \texttt{MAGE\_K=8192} is inactive as a support
constraint for V2. Experimental target sparsity and realized physical work
must be reported separately from $\tau_V$.
\end{minipage}
\end{table*}
% END EMBEDDED METHODS SECTION
\end{document}
